\chapter{冒险旋钮}
% 完整版：chapters/08_nora.tex

你在挑地方吃午饭。去熟悉的那家小馆，还是试一家新的？如果总去熟悉的，就永远不会知道拐角处还有更好的。如果总在尝新，就会常常吃得很糟。平衡点在哪里？这个问题——继续探索，还是利用已经找到的——摆在每一个会学习的生物面前。在大脑里，负责它的看来是\nt{NORA}电台——去甲肾上腺素。

\section*{对比度旋钮}

去甲肾上腺素改变的是神经元对信号的反应有多陡峭。作用弱时，神经元像一个平滑的调光器：信号大一点，响应就亮一点。作用强时，它像一个开关：阈值以下——沉默，阈值以上——全力响应。更准确地说，实验表明，去甲肾上腺素压制神经元背景“闲聊”的程度，要大于压制它们对真实信号的响应；“对比度旋钮”是描述这一效应的便利模型。

\pencilfig{8}{\pencilart{8}}{冒险旋钮：往左拧——机器尝试新东西，往右拧——果断选择熟悉的。}

\section*{思考还是决定}

回想第三章里那两群互相抑制的神经元。对比度低时，两群都在工作，弱的一群稍微安静一些：网络还在“思考”，在权衡各个选项。对比度一旦超过某个阈值，网络就猛然切换到“决定了”的模式：一群胜出，另一群沉默。一个旋钮就能让机器在斟酌与决断之间切换。

\section*{探索多少}

加入噪声。对比度低时，噪声很容易压过选项之间的微小差别，机器常常尝试不是最好的那个选项——它在探索。对比度高时，噪声左右不了什么，机器总是选它认为最好的——它在利用。

在我们的模型里，世界时不时会变：最好的选项变成最差的。收益随旋钮的变化呈倒 U 形。对比度太低——机器盲目游荡。太高——它察觉不到世界已经变了，固执地往那家已经变差的小馆跑。世界变得越频繁，最佳设置就越低：在多变的世界里，多探索更划算。

\pencilfig{108}{%
% points: models/nora_explore.py, reward as a share of the best possible, gain 0.5 ... 64 (doubling); the y axis starts at 0.70, hence the break mark
\draw[pen] (0,0) -- (6.3,0); \draw[pen] (0,0) -- (0,0.18); \draw[pen] (0,0.34) -- (0,2.4);
\draw[pcGray, line width=0.6pt] (-0.1,0.14) -- (0.1,0.22); \draw[pcGray, line width=0.6pt] (-0.1,0.3) -- (0.1,0.38);
\node[lbl, anchor=north] at (3.0,-0.05) {对比度旋钮};
\node[lbl, anchor=south, rotate=90] at (-0.1,1.3) {收益};
\draw[penlite=pcBlue] plot[smooth] coordinates {(0.00,0.55) (0.85,0.79) (1.70,1.19) (2.55,1.68) (3.40,2.00) (4.25,2.07) (5.10,2.01) (5.95,1.91)};
\draw[penlite=pcRed] plot[smooth] coordinates {(0.00,0.55) (0.85,0.69) (1.70,0.93) (2.55,1.24) (3.40,1.40) (4.25,1.28) (5.10,1.06) (5.95,0.89)};
\draw[dotpen=pcBlue, wash=pcBlue] (4.25,2.07) circle (0.07);
\draw[dotpen=pcRed, wash=pcRed] (3.40,1.40) circle (0.07);
\node[lbl, text=pcBlue, anchor=south] at (4.25,2.15) {世界很少变化};
\node[lbl, text=pcRed, anchor=north] at (3.40,1.30) {经常变化};
\node[lbl, anchor=north west, align=left] at (0.05,-0.35) {盲目\\游荡}; \node[lbl, anchor=north east, align=right] at (6.3,-0.35) {察觉不到\\变化};
}{旋钮处于不同位置时模型的收益。每条曲线都有一个峰；世界变化越频繁，峰就越低、越靠左。}

\section*{谁在拧旋钮}

变化的天然标志是意外：误差比平常大了。按照一种假说，去甲肾上腺素报告的正是这种意外的变化。这里有个微妙之处：持续的背景水平的去甲肾上腺素更多地与探索和分心有关，而针对重要事件的短暂爆发则与专注和决断有关。爆发也许像计算机里的中断一样工作：“放下手头的事，世界变了，重新制定计划。”

这套系统有一个从外面就看得见的间接标志：当某件事意外改变、人开始学得更快时，瞳孔会放大。不过瞳孔跟踪的不只是去甲肾上腺素，所以它只是一条线索，而不是测量仪器。

还要坦白一点：在我们的模型里，根据意外来调节旋钮，比起最佳的固定设置只多赢了一点点。这种调节究竟在什么情况下才划得来，还有待弄清。

\fullversion{8}
